% methodology.tex -- Sequencer Paper (BNR-2026-009)

\section{Methodology}

\subsection{Research Design}

The investigation follows a three-phase approach:

\textbf{Phase 1: Literature Survey.} Systematic analysis of eight production L2 sequencer
architectures to identify universal patterns, forced inclusion mechanisms, and published
performance benchmarks. Sources include protocol documentation, source code repositories,
and technical blog posts from each project.

\textbf{Phase 2: Prototype Implementation.} A Go prototype implementing the core
sequencer components: FIFO mempool, forced inclusion queue, and tick-based block
production. The prototype omits EVM execution (which is the domain of RU-L1,
BNR-2026-008~\cite{bnr2026008evm}) to isolate sequencer-specific performance
characteristics.

\textbf{Phase 3: Benchmark Evaluation.} Comprehensive benchmarks across four dimensions:
(a)~block production scaling with increasing transaction counts, (b)~forced inclusion
under cooperative and adversarial conditions, (c)~concurrent mempool access with multiple
producer goroutines, and (d)~high-throughput direct measurement without ticker jitter.

\subsection{Prototype Architecture}

The prototype consists of four Go source files (${\sim}$500~LOC total):

\begin{table}[htbp]
\caption{Prototype Source Code Structure}
\label{tab:prototype-structure}
\centering
\begin{tabular}{llr}
\toprule
\textbf{File} & \textbf{Component} & \textbf{LOC} \\
\midrule
\texttt{types.go} & Transaction, Block, Config, Metrics & 178 \\
\texttt{mempool.go} & FIFO Mempool (Add, Drain, Capacity) & 139 \\
\texttt{forced\_inclusion.go} & Forced Inclusion Queue (Arbitrum-style) & 115 \\
\texttt{sequencer.go} & Block Production, Run Loop, Metrics & 230 \\
\midrule
\textbf{Total} & & \textbf{662} \\
\bottomrule
\end{tabular}
\end{table}

\textbf{Mempool} (Listing~\ref{lst:mempool}): Thread-safe FIFO queue backed by a
\texttt{sync.Mutex}-protected slice. Each transaction receives a monotonically
increasing sequence number via atomic counter. The \texttt{Drain(n)} operation extracts
up to $n$ transactions from the front, preserving FIFO order.

\begin{lstlisting}[language=Go, caption={FIFO Mempool core operations}, label={lst:mempool}]
func (m *Mempool) Add(tx *Transaction) bool {
  m.mu.Lock()
  defer m.mu.Unlock()
  if len(m.txs) >= m.maxSize {
    return false // capacity enforcement
  }
  tx.Sequence = atomic.AddUint64(
    &m.seqCounter, 1)
  m.txs = append(m.txs, tx)
  return true
}

func (m *Mempool) Drain(n int) []*Transaction {
  m.mu.Lock()
  defer m.mu.Unlock()
  take := min(n, len(m.txs))
  result := m.txs[:take]
  m.txs = m.txs[take:]
  return result
}
\end{lstlisting}

\textbf{Forced Inclusion Queue}: Maintains a FIFO queue of L1-submitted transactions,
each with a submission timestamp and configurable deadline. The \texttt{DrainDue(now)}
method returns all transactions whose deadline has expired, preserving submission order.

\textbf{Block Producer}: The \texttt{ProduceBlock()} method is the core benchmarked
function. It drains forced transactions first, then mempool transactions up to the
per-block limit, assembles the block with sequential numbering and parent hash linkage,
and updates metrics.

\subsection{Benchmark Methodology}

\textbf{Environment}: Windows 11 Pro, 20 cores (Intel 3.30~GHz), Go 1.22.10.

\textbf{Repetitions}: 30 repetitions per scaling configuration with different seeds.
Scenario tests use fixed configurations to ensure reproducibility.

\textbf{Direct measurement}: High-throughput benchmarks pre-load transactions into the
mempool before timing block production, eliminating ticker jitter and measuring pure
sequencer overhead.

\textbf{Go native benchmarks}: \texttt{go test -bench} provides nanosecond-precision
measurements of individual operations (mempool insert, mempool drain, block production)
with automatic iteration count calibration.

\textbf{Metrics}: All benchmarks report:
\begin{itemize}
  \item Block production latency (microseconds): time to execute \texttt{ProduceBlock()}
  \item Mempool insert throughput (tx/s): insertion rate under load
  \item FIFO accuracy (\%): percentage of transactions in correct arrival order
  \item Forced inclusion rate (\%): percentage of forced transactions included within deadline
  \item Forced inclusion latency (ms): time from submission to block inclusion
\end{itemize}

\subsection{Benchmark Validation Targets}

Table~\ref{tab:validation-targets} presents the expected performance ranges derived
from published benchmarks, against which our measurements are validated.

\begin{table}[htbp]
\caption{Benchmark Validation Targets}
\label{tab:validation-targets}
\centering
\begin{tabular}{lrl}
\toprule
\textbf{Metric} & \textbf{Expected} & \textbf{Source} \\
\midrule
Block production & $<$50~ms at 100~TPS & Geth, OP Stack \\
Mempool insert & $>$10{,}000~tx/s & Geth txpool \\
Forced inclusion detect. & $<$100~ms/L1 block & ethers.js \\
Block sealing overhead & $<$5~ms & Consensus layer \\
FIFO ordering accuracy & 100\% & By design \\
\bottomrule
\end{tabular}
\end{table}

\subsection{Test Suite}

The prototype includes 8 unit tests and 6 scenario tests:

\textbf{Unit tests}: Mempool FIFO ordering, capacity enforcement (insert rejected when
full), forced inclusion deadline expiry, block production with mixed forced and mempool
transactions, empty block production, sequence number monotonicity, concurrent mempool
access safety, and metrics accuracy.

\textbf{Scenario tests}: Steady-state 100~TPS, steady-state 500~TPS, burst 1{,}000~TPS,
maximum capacity stress, cooperative forced inclusion (20~forced transactions at 1-second
deadline), and adversarial censoring (forced inclusion under simulated sequencer delay).
